\chapter{Continuous Distributions and the Poisson Process}

\paragraph{Exercise 8.1} Let $X$ and $Y$ be independent, uniform random variables
on $[0,1]$. We want to determine the density and the distribution function of the
sum of the two uniform random variables $X$ and $Y$. Let $Z$ be a random variable,
such that $Z = X + Y$. Hence, we want to determine $f_Z$ and $F_Z$. For all
$z \in \mathbb{R}$,
\[
    F_Z(z) = \int_{- \infty}^{z} f_Z(u) \, du,
\]
and
\[
    f_Z(z)
        = \int_{- \infty}^{\infty} f_{X,Y}(x,z-x) \, dx
        = \int_{- \infty}^{\infty} f_X(x)\cdot f_Y(z-x) \, dx.
\]
Furthermore, $f_Z(z) = 0$ for $z < 0$, and also for $z \geq 2$. In regard to the
case $z \in [0,2]$, we will distinguish two cases.
\begin{enumerate}
    \item[(i)] Assume that $z \in [0,1]$. In order to have $f_Y(z-x) = 1$, we need
    $z-x \geq 0$, that is, $x \leq z$. Consequently, it suffices to integrate from
    $x=0$ to $x = z$,
    \[
        f_z(z)
            = \int_{0}^{z} f_X(x)\cdot f_Y(z-x) \, dx
            = \int_{0}^{z} 1 \, dx
            = z.
    \]

    \item[(ii)] Assume that $z \in [1,2]$. In order to have $f_Y(z-x) = 1$, we
    need $z-x \leq 1$, that is, we need $x \geq z-1$. So in this case, we integrate
    from $z-1$ to 1,
    \[
    f_z(z)
        = \int_{z-1}^{1} f_X(x)\cdot f_Y(z-x) \, dx
        = \int_{z-1}^{1} 1 \, dx
        = \left[ x \right]_{z-1}^1
        = 1 - (z-1)
        = 2 - z.
    \]
\end{enumerate}
Thus, the density function of $Z$ is defined as
\[
    f_Z: z \mapsto
    \begin{cases}
        z, &\text{if }z  \in [0,1]; \\
        2-z, &\text{if }z \in [1,2]; \\
        0,  &\text{else.}
    \end{cases}
\]
Let us now consider the distribution function of $Z$. $F_Z(z) = 0$ for $z < 0$,
and $F_Z(z) = 1$ for $z \geq 2$. In regard to the case $z \in [0,2]$, we will
examine the two cases $z \in [0,1]$ and $z \in [1,2]$.
\begin{enumerate}
    \item[(i)] Assume that $z \in [0,1]$.
    \[
        F_Z(z)
            = \int_{- \infty}^{z} f_Z(u) \, du
            = \int_{0}^{z} u \, du
            = \left[\frac{u^2}{2}\right]_0^z
            = \frac{z^2}{2}.
    \]

    \item[(ii)] Assume that $z \in [1,2]$.
    \begin{align*}
        F_Z(z)
            &= \int_{- \infty}^{z} f_Z(u) \, du
            = \int_{0}^{1} u du + \int_{1}^{z} 2-u \, du \\
            &= \left[\frac{u^2}{2}\right]_0^1 + \left[2u -  \frac{u^2}{2}\right]_1^z
            = \frac{1}{2} + 2z - \frac{z^2}{2} - 2 + \frac{1}{2}
            = - \frac{z^2}{2} + 2z - 1.
    \end{align*}
\end{enumerate}
Thus, the distribution function of $Z$ is given by
\[
    F_Z: z \mapsto
    \begin{cases}
        0,                          &\text{if }z  < 0; \\
        \frac{z^2}{2},              &\text{if }z \in [0,1]; \\
        - \frac{z^2}{2} + 2z - 1,   &\text{if }z \in [1,2]; \\
        1,                          &\text{if }z > 2.
    \end{cases}
\]

\paragraph{Exercise 8.2} Let $X$ and $Y$ be independent, exponentially distributed
random variables with parameter 1. Let $Z$ be a random variable,
such that $Z = X + Y$. For all $z \in \mathbb{R}$,
\[
    F_Z(z) = \int_{- \infty}^{z} f_Z(u) \, du,
\]
and
\[
    f_Z(z)
        = \int_{- \infty}^{\infty} f_{X,Y}(x,z-x) \, dx
        = \int_{- \infty}^{\infty} f_X(x)\cdot f_Y(z-x) \, dx.
\]
Clearly, $f_Z(z) = 0$ for $z < 0$. Consider the case $z \geq 0$. In order to have
$f_Y(z-x) \not= 0$, we need $z-x \geq 0$, that is, $x \leq z$. Consequently, it
suffices to integrate from $x=0$ to $x = z$,
\[
    f_z(z)
        = \int_{0}^{z} f_X(x)\cdot f_Y(z-x) \, dx
        = \int_{0}^{z} e^{-x} \cdot e^{-(z-x)} \, dx
        = \int_{0}^{z} e^{-x-(z-x)} \, dx
        = e^{-z} \int_{0}^{z} 1 \, dx
        = z \cdot e^{-z}.
\]
Thus, the density function of $Z$ is defined as
\[
    f_Z: z \mapsto
    \begin{cases}
        z \cdot e^{-z}, &\text{if }z  \geq 0; \\
        0,  &\text{else.}
    \end{cases}
\]
Let us now consider the distribution function of $Z$. $F_Z(z) = 0$ for $z < 0$.
In regard to the case $z \geq 0$, one has
\[
    F_Z(z)
        = \int_{- \infty}^{z} f_Z(u) \, du
        = \int_{0}^{z} z \cdot e^{-z} \, du
        = z \cdot e^{-z} \int_{0}^{z} 1 \, du
        = z^2 \cdot e^{-z}.
\]
Thus, the distribution function of $Z$ is given by
\[
    F_Z: z \mapsto
    \begin{cases}
        z^2 \cdot e^{-z},   &\text{if }z \geq 0; \\
        0,                  &\text{else.}
    \end{cases}
\]

\paragraph{Exercise 8.3} Let $X$ be a uniform random variable on [0,1]. Then,
\[
    \pr\left(X \leq \frac{1}{2} \mid \frac{1}{4} \leq X \leq \frac{3}{4}\right)
    = \frac{\pr\left(\frac{1}{4} \leq X \leq \frac{1}{2}\right)}{\pr\left(\frac{1}{4} \leq X \leq \frac{3}{4}\right)}
    = \frac{\frac{1}{2} - \frac{1}{4}}{\frac{3}{4} - \frac{1}{4}}
    = \frac{1}{2},
\]
and
\begin{align*}
    \pr\left(X \leq \frac{1}{4} \mid \left(X \leq \frac{1}{3}\right) \cup \left( X \geq \frac{2}{3}\right) \right)
    &= \frac{\pr\left(\left(X \leq \frac{1}{4} \right) \cap \left( \left(X \leq \frac{1}{3}\right) \cup \left( X \geq \frac{2}{3}\right) \right) \right)}{\pr\left(\left(X \leq \frac{1}{3}\right) \cup \left( X \geq \frac{2}{3}\right) \right)} \\
    &= \frac{\pr\left(X \leq \frac{1}{4}\right)}{\pr\left(X \leq \frac{1}{3} \right) + \pr\left(X \geq \frac{2}{3}\right)} \\
    & = \frac{\frac{1}{4}}{\frac{1}{3} + \frac{1}{3}} \\
    &= \frac{3}{8}.
\end{align*}
